% Subseção da F-MET-OPC (comunidades, RF-MET-06).
% Fontes: docs/plano-de-desenvolvimento/features/feature-F-MET-OPC.md.
% Números do feed real: python benchmarks/medir_metricas_extensoes.py
% dados/brutos/bhtrans.zip --repeticoes 3 (medição final pendente).
\subsection{Comunidades}\label{sec:F-MET-OPC}

A aba Comunidades da página Centralidades agrupa as paradas mais ligadas entre si do que
com o resto da rede, por propagação de rótulos sobre $G_P$ tratado como não dirigido e sem
pesos. Cada parada começa com o próprio rótulo; a cada iteração, numa ordem embaralhada com
semente fixa, adota o rótulo mais frequente entre as vizinhas, mantendo o atual em caso de
empate com ele e, senão, o menor. Para quando nada muda. Cada iteração custa $O(V + E)$. A
qualidade da partição é a modularidade de Newman, $Q = \sum_c [L_c/m - (d_c/2m)^2]$,
implementada por nós e conferida com a do NetworkX em 15 grafos aleatórios; nos testes, a
propagação recupera cliques ligados por uma ponte. No feed real, com a semente 2026, a
propagação convergiu em 11 iterações (\pendente{tempo}) para 604 comunidades, a maior com
141 paradas, e modularidade $0{,}806$, igual à do NetworkX para a mesma partição.
\emph{Limitação:} o resultado depende da semente (com cinco sementes, de 587 a 608
comunidades e $Q$ entre $0{,}806$ e $0{,}815$) e ignora o tempo das ligações; as comunidades
saem pequenas (a maior tem 1,5\% das paradas), sem formar regiões da cidade.
